\documentclass[11pt,a4paper]{article}
\usepackage{geometry}
\geometry{top=2.5cm,bottom=2.5cm,left=2.3cm,right=2.3cm}

\usepackage{xeCJK}
\setCJKmainfont{PingFang SC}[BoldFont=PingFang SC Semibold]
\setCJKsansfont{PingFang SC}
\setCJKmonofont{PingFang SC}

\setmainfont{Helvetica Neue}[BoldFont=Helvetica Neue Bold]
\setsansfont{Helvetica Neue}
\setmonofont{Menlo}[Scale=0.88]

\usepackage{xcolor}
\definecolor{navyblue}{HTML}{0F172A}
\definecolor{accentblue}{HTML}{0284C7}
\definecolor{borderblue}{HTML}{38BDF8}
\definecolor{bgblue}{HTML}{F0F9FF}

\definecolor{warnamber}{HTML}{D97706}
\definecolor{borderamber}{HTML}{F59E0B}
\definecolor{bgamber}{HTML}{FFFBEB}

\definecolor{dangerred}{HTML}{DC2626}
\definecolor{borderred}{HTML}{EF4444}
\definecolor{bgred}{HTML}{FEF2F2}

\definecolor{tipgreen}{HTML}{16A34A}
\definecolor{bordergreen}{HTML}{22C55E}
\definecolor{bggreen}{HTML}{F0FDF4}

\definecolor{noteslate}{HTML}{475569}
\definecolor{borderslate}{HTML}{94A3B8}
\definecolor{bgslate}{HTML}{F8FAFC}

\definecolor{codebg}{HTML}{F1F5F9}
\definecolor{codeframe}{HTML}{CBD5E1}

\usepackage{tcolorbox}
\tcbuselibrary{skins,breakable}

% Custom GFM Alert Boxes
\newtcolorbox{gfmImportant}[1][]{
  colback=bgblue,
  colframe=accentblue,
  leftrule=4pt,
  rightrule=0.5pt,
  toprule=0.5pt,
  bottomrule=0.5pt,
  arc=3pt,
  boxsep=4pt,
  left=8pt,right=8pt,top=6pt,bottom=6pt,
  breakable,
  title={\textbf{\color{navyblue} IMPORTANT #1}},
  coltitle=navyblue
}

\newtcolorbox{gfmWarning}[1][]{
  colback=bgamber,
  colframe=warnamber,
  leftrule=4pt,
  rightrule=0.5pt,
  toprule=0.5pt,
  bottomrule=0.5pt,
  arc=3pt,
  boxsep=4pt,
  left=8pt,right=8pt,top=6pt,bottom=6pt,
  breakable,
  title={\textbf{\color{warnamber} WARNING #1}},
  coltitle=warnamber
}

\newtcolorbox{gfmNote}[1][]{
  colback=bgslate,
  colframe=noteslate,
  leftrule=4pt,
  rightrule=0.5pt,
  toprule=0.5pt,
  bottomrule=0.5pt,
  arc=3pt,
  boxsep=4pt,
  left=8pt,right=8pt,top=6pt,bottom=6pt,
  breakable,
  title={\textbf{\color{noteslate} NOTE #1}},
  coltitle=noteslate
}

\newtcolorbox{gfmTip}[1][]{
  colback=bggreen,
  colframe=tipgreen,
  leftrule=4pt,
  rightrule=0.5pt,
  toprule=0.5pt,
  bottomrule=0.5pt,
  arc=3pt,
  boxsep=4pt,
  left=8pt,right=8pt,top=6pt,bottom=6pt,
  breakable,
  title={\textbf{\color{tipgreen} TIP #1}},
  coltitle=tipgreen
}

\newtcolorbox{gfmCaution}[1][]{
  colback=bgred,
  colframe=dangerred,
  leftrule=4pt,
  rightrule=0.5pt,
  toprule=0.5pt,
  bottomrule=0.5pt,
  arc=3pt,
  boxsep=4pt,
  left=8pt,right=8pt,top=6pt,bottom=6pt,
  breakable,
  title={\textbf{\color{dangerred} CAUTION #1}},
  coltitle=dangerred
}

\newtcolorbox{gfmBlockquote}[1][]{
  colback=bgslate,
  colframe=borderslate,
  leftrule=3pt,
  rightrule=0.2pt,
  toprule=0.2pt,
  bottomrule=0.2pt,
  arc=2pt,
  boxsep=3pt,
  left=6pt,right=6pt,top=4pt,bottom=4pt,
  breakable
}

\usepackage{tabularx}
\usepackage{booktabs}
\usepackage{colortbl}
\usepackage{multirow}

\usepackage{listings}
\lstset{
  basicstyle=\ttfamily\small,
  backgroundcolor=\color{codebg},
  frame=single,
  rulecolor=\color{codeframe},
  breaklines=true,
  breakatwhitespace=true,
  showstringspaces=false,
  tabsize=2,
  captionpos=b,
  keywordstyle=\color{accentblue}\bfseries,
  commentstyle=\color{noteslate}\itshape,
  stringstyle=\color{tipgreen}
}

\usepackage{graphicx}
\usepackage{fancyhdr}
\usepackage{lastpage}
\usepackage{hyperref}
\hypersetup{
  colorlinks=true,
  linkcolor=accentblue,
  urlcolor=accentblue,
  citecolor=accentblue,
  pdfborder={0 0 0}
}

\pagestyle{fancy}
\fancyhf{}
\fancyhead[L]{\small\color{noteslate}Atlas GitOps Control Plane Design}
\fancyhead[R]{\small\color{noteslate}Atlas Architecture Specification}
\fancyfoot[C]{\small\color{noteslate}Page \thepage\ of \pageref{LastPage}}
\renewcommand{\headrulewidth}{0.4pt}
\renewcommand{\footrulewidth}{0.4pt}

\title{\textbf{\Large Atlas GitOps Control Plane Design}}
\author{Atlas Architecture Group}
\date{\today}

\begin{document}
\maketitle
\tableofcontents
\newpage

\textbf{Version:} v1.0.3

\textbf{Status:} Frozen Production Baseline (Two-Level DAG Architecture Freeze Point)

\textbf{Domain:} GitOps Control Plane / Platform Engineering

\vspace{0.5em}\hrule\vspace{0.5em}
\subsection{0. 文档职责与规范语义}
本文档是：

\begin{gfmBlockquote}
\textbf{Atlas Architecture Design v1.0.2}
\end{gfmBlockquote}
在 GitOps / Argo CD 领域中的规范性执行设计。

本文档负责定义：

\begin{itemize}
  \item External Root Anchor；
  \item AppProject Trust Chain；
  \item Root Macro DAG；
  \item Platform Capability DAG；
  \item Workload Control Model；
  \item Sync Wave；
  \item Application Health Gate；
  \item Environment Hydration；
  \item Child Automated Sync；
  \item Deletion Protection；
  \item Bootstrap Adoption；
  \item Break-Glass Recovery。
\end{itemize}
本文不得改变 Architecture Design 中定义的平台级不变量。

\vspace{0.5em}\hrule\vspace{0.5em}
\subsection{1. GitOps Ownership \& Trust Boundary}
Atlas GitOps 控制面采用四阶控制权模型。

\begin{lstlisting}[language=text]
Git
 │ Definition
 ▼
Bootstrap
 │ Instantiation
 ▼
Argo CD
 │ Desired-State Reconciliation
 ▼
Kubernetes Controllers / Operators
   Runtime Behavior
\end{lstlisting}
\vspace{0.5em}\hrule\vspace{0.5em}
\subsubsection{1.1 Git owns Definition}
Git 保存：

\begin{itemize}
  \item Root Anchor Template；
  \item Application Graph；
  \item AppProject；
  \item Platform Desired State；
  \item Workload Desired State；
  \item Policy；
  \item Overlay；
  \item Version Lock。
\end{itemize}
Git 是 Desired State Definition Authority。

\vspace{0.5em}\hrule\vspace{0.5em}
\subsubsection{1.2 Bootstrap owns Instantiation}
Bootstrap 负责：

\begin{enumerate}
  \item 安装最小 Argo CD Seed；
  \item 创建 \_\_INLINE\_CODE\_0\_\_；
  \item 注入与 Full Desired State 相同的 Application Health Capability；
  \item 根据环境上下文渲染 Root Anchor；
  \item 将 External Root Anchor 注入 Kubernetes。
\end{enumerate}
Bootstrap 完成后退出。

\vspace{0.5em}\hrule\vspace{0.5em}
\subsubsection{1.3 Argo CD owns GitOps Reconciliation}
Argo CD 对所有进入 GitOps Domain 的资源负责：

\begin{itemize}
  \item Sync；
  \item Drift Detection；
  \item Reconciliation；
  \item Health；
  \item Prune。
\end{itemize}
External Root Anchor 本身除外。

External Root Anchor 不存在 Parent Application，因此不属于 parent-driven steady-state reconciliation。

\vspace{0.5em}\hrule\vspace{0.5em}
\subsubsection{1.4 Operators own Runtime Behavior}
例如：

\begin{lstlisting}[language=text]
Argo CD
 ↓
FlinkDeployment
 ↓
Flink Operator
 ↓
Runtime
\end{lstlisting}
Argo CD 不尝试取代 Domain Operator。

\vspace{0.5em}\hrule\vspace{0.5em}
\subsection{2. Canonical AppProject Model}
Atlas 正式冻结三个 Canonical AppProject Identifier。

\begin{lstlisting}[language=text]
atlas-bootstrap
platform-project
workload-project
\end{lstlisting}
任何实现、图示、Manifest 与 Conformance Audit MUST 使用上述准确名称。

禁止使用以下模糊别名代替实际资源名称：

\begin{lstlisting}[language=text]
platform
workload
bootstrap-project
\end{lstlisting}
\vspace{0.5em}\hrule\vspace{0.5em}
\subsection{3. Three-Tier Trust Isolation}
\vspace{0.5em}\hrule\vspace{0.5em}
\subsubsection{3.1 Tier-0 — \_\_INLINE\_CODE\_0\_\_}
\_\_INLINE\_CODE\_0\_\_ 由 Bootstrap 命令式创建。

其唯一职责是承载完成 Trust Chain 自举所必需的 Application。

允许：

\begin{itemize}
  \item External Root Anchor；
  \item Project Bootstrap。
\end{itemize}
禁止：

\begin{itemize}
  \item 普通 Platform Application；
  \item Workload Application；
  \item Developer Workload；
  \item 日常平台组件。
\end{itemize}
\_\_INLINE\_CODE\_0\_\_ 是 Bootstrap Capability Sandbox，不是普通租户 Project。

\vspace{0.5em}\hrule\vspace{0.5em}
\subsubsection{3.2 Tier-1 — \_\_INLINE\_CODE\_0\_\_}
由 Project Bootstrap 声明式创建。

负责：

\begin{itemize}
  \item Platform Control；
  \item Argo CD Self-Management；
  \item Infrastructure；
  \item Governance；
  \item Operator；
  \item Controller；
  \item Workload Orchestrator；
  \item 需要高权限的管理组件。
\end{itemize}
权限必须采用显式 allow-list，而不是无约束开放。

\vspace{0.5em}\hrule\vspace{0.5em}
\subsubsection{3.3 Tier-2 — \_\_INLINE\_CODE\_0\_\_}
由 Project Bootstrap 声明式创建。

用于最终业务工作负载。

其原则是：

\begin{gfmBlockquote}
Restricted by default.
\end{gfmBlockquote}
必须禁止或严格限制：

\begin{itemize}
  \item CRD；
  \item ClusterRole；
  \item ClusterRoleBinding；
  \item Cluster-wide privileged resources；
  \item 任意 Platform Namespace；
  \item Argo CD Control Namespace。
\end{itemize}
Namespaced：

\begin{itemize}
  \item ServiceAccount；
  \item Role；
  \item RoleBinding；
\end{itemize}
按平台策略受控开放。

\vspace{0.5em}\hrule\vspace{0.5em}
\subsection{4. Bootstrap Trust Chain}
Atlas 的完整无环 Bootstrap Chain 为：

\begin{lstlisting}[language=text]
                 Bootstrap
                     │
          ┌──────────┴──────────┐
          ▼                     ▼
     Argo CD Seed       atlas-bootstrap
                          AppProject
                               │
                               ▼
                  External Root Anchor
                  project: atlas-bootstrap
                               │
                               ▼
                 Project Bootstrap App
                      wave: -110
                  project: atlas-bootstrap
                               │
                ┌──────────────┴──────────────┐
                ▼                             ▼
        platform-project               workload-project
                │                             │
                └──────────────┬──────────────┘
                               ▼
                    Normal GitOps Control
\end{lstlisting}
不存在：

\begin{lstlisting}[language=text]
Application requires Project
        ↑
Project requires same Application
\end{lstlisting}
的循环依赖。

\vspace{0.5em}\hrule\vspace{0.5em}
\subsection{5. External Root Anchor}
External Root Anchor 是整个 GitOps Control Graph 的外部根节点。

其模板位于：

\begin{lstlisting}[language=text]
bootstrap/argocd/root-app.yaml.tpl
\end{lstlisting}
Bootstrap 根据：

\begin{lstlisting}[language=text]
env/
\end{lstlisting}
中的环境上下文进行渲染，并将实例写入：

\begin{lstlisting}[language=text]
.state/root-app.yaml
\end{lstlisting}
随后注入 Kubernetes。

其 Desired Definition 来源于 Git，但该活体对象不存在 Parent Application。

因此：

\begin{gfmBlockquote}
Git owns its contract, Bootstrap owns its instantiation, but no parent Application owns its reconciliation.
\end{gfmBlockquote}
这是有意设计，不属于配置缺陷。

\vspace{0.5em}\hrule\vspace{0.5em}
\subsubsection{5.1 Root Anchor Security Consequence}
由于 Root Anchor 不被父级 Argo CD Application 修复：

\begin{lstlisting}[language=text]
kubectl edit Application <root>
\end{lstlisting}
理论上可以导致：

\begin{lstlisting}[language=text]
Live Root
≠
Git Template
\end{lstlisting}
因此 External Root Anchor 必须通过其他控制机制保护：

\begin{itemize}
  \item Kubernetes RBAC；
  \item Argo RBAC；
  \item Admission Policy；
  \item Audit；
  \item Tier-0 Human Judgment；
  \item Drift Detection。
\end{itemize}
未来 Kyverno / 等价准入治理应成为该 Trust Boundary 的重要补强。

\vspace{0.5em}\hrule\vspace{0.5em}
\subsection{6. Two-Level Reconciliation DAG}
Atlas v1.0.3 正式放弃：

\begin{gfmBlockquote}
Flat Global Root Component DAG
\end{gfmBlockquote}
同时禁止：

\begin{gfmBlockquote}
Filesystem-shaped Deep Nested DAG.
\end{gfmBlockquote}
正式模型为：

\begin{gfmBlockquote}
\textbf{Minimal Root Macro DAG + Flat Capability DAG}
\end{gfmBlockquote}
\vspace{0.5em}\hrule\vspace{0.5em}
\subsection{7. Level 1 — Root Macro DAG}
\_\_INLINE\_CODE\_0\_\_ 只负责 Trust Transition。

标准结构：

\begin{lstlisting}[language=text]
gitops/root/
├── base/
│   ├── kustomization.yaml
│   ├── project-bootstrap-app.yaml
│   ├── platform-control-app.yaml
│   └── workload-control-app.yaml
│
└── overlays/
    ├── local-orbstack/
    │   └── kustomization.yaml
    └── prod/
        └── kustomization.yaml
\end{lstlisting}
Root SHALL NOT 直接列举：

\begin{lstlisting}[language=text]
flink
envoy
prometheus
redpanda
minio
loki
tempo
...
\end{lstlisting}
\vspace{0.5em}\hrule\vspace{0.5em}
\subsubsection{7.1 Root Macro Waves}
Root Synchronization Scope 定义：

\begin{lstlisting}[language=text]
-110  Project Bootstrap
-100  Platform Control
   0  Workload Control
\end{lstlisting}
控制流程：

\begin{lstlisting}[language=text]
Project Bootstrap Healthy
          ↓
Platform Control starts
          ↓
Platform Control Healthy
          ↓
Workload Control starts
\end{lstlisting}
Root 的 Wave 与 Platform Control 内部 Wave 属于不同 Sync Scope。

数字相同不表示同一全局时间点。

\vspace{0.5em}\hrule\vspace{0.5em}
\subsection{8. Level 2 — Platform Capability DAG}
Platform Control Application 指向：

\begin{lstlisting}[language=text]
gitops/platform/applications/overlays/<environment>/
\end{lstlisting}
例如：

\begin{lstlisting}[language=text]
gitops/platform/applications/
├── base/
│   ├── kustomization.yaml
│   ├── foundation-app.yaml
│   ├── argocd-self-app.yaml
│   ├── sealed-secrets-app.yaml
│   ├── cert-manager-app.yaml
│   ├── flink-operator-app.yaml
│   ├── envoy-gateway-app.yaml
│   ├── kube-prometheus-stack-app.yaml
│   ├── redpanda-app.yaml
│   └── minio-app.yaml
│
└── overlays/
    ├── local-orbstack/
    └── prod/
\end{lstlisting}
这里保存的是：

\begin{gfmBlockquote}
Application Graph Metadata
\end{gfmBlockquote}
而不是组件实现。

\vspace{0.5em}\hrule\vspace{0.5em}
\subsection{9. Platform Capability Wave DAG}
Platform Control Synchronization Scope 内采用：

\begin{lstlisting}[language=text]
-100 Foundation
 -90 Management Plane
 -50 Operators
 -30 Infrastructure Controllers
 -10 Platform Services
\end{lstlisting}
逻辑如下：

\begin{lstlisting}[language=text]
                    -100
                 Foundation
                     │
                     ▼
                     -90
               Management Plane
                     │
                     ▼
                     -50
                  Operators
                     │
                     ▼
                     -30
           Infrastructure Controllers
                     │
                     ▼
                     -10
              Platform Services
\end{lstlisting}
\vspace{0.5em}\hrule\vspace{0.5em}
\subsubsection{9.1 Wave -100 — Foundation}
包括：

\begin{itemize}
  \item Namespaces；
  \item ResourceQuotas；
  \item LimitRanges；
  \item 其他 Kubernetes 原生逻辑边界。
\end{itemize}
环境强绑定资源必须通过 Environment Hydration 进入。

\vspace{0.5em}\hrule\vspace{0.5em}
\subsubsection{9.2 Wave -90 — Management Plane}
包括：

\begin{itemize}
  \item \_\_INLINE\_CODE\_0\_\_；
  \item Sealed Secrets；
  \item cert-manager；
  \item OIDC；
  \item Policy Engine；
  \item ESO；
  \item 其他平台自身管理与治理能力。
\end{itemize}
其中 Secret Management Controller 必须早于其 Consumer。

\vspace{0.5em}\hrule\vspace{0.5em}
\subsubsection{9.3 Wave -50 — Operators}
包括：

\begin{itemize}
  \item Flink Kubernetes Operator；
  \item Spark Operator；
  \item 其他扩展 Kubernetes API 的 Domain Operator。
\end{itemize}
必须在对应 CR 出现前确保：

\begin{itemize}
  \item CRD 已注册；
  \item Controller Ready；
  \item Webhook Ready；
  \item Application Healthy。
\end{itemize}
\vspace{0.5em}\hrule\vspace{0.5em}
\subsubsection{9.4 Wave -30 — Infrastructure Controllers}
包括：

\begin{itemize}
  \item Envoy Gateway Controller；
  \item Prometheus Operator；
  \item 其他底层 Controller。
\end{itemize}
“Controller”是依赖角色，不是目录领域。

例如：

\begin{lstlisting}[language=text]
Envoy Gateway
\end{lstlisting}
仍然属于：

\begin{lstlisting}[language=text]
platform/networking/
\end{lstlisting}
而不是：

\begin{lstlisting}[language=text]
platform/controllers/
\end{lstlisting}
\vspace{0.5em}\hrule\vspace{0.5em}
\subsubsection{9.5 Wave -10 — Platform Services}
包括：

\begin{itemize}
  \item Redpanda / Kafka；
  \item MinIO；
  \item Database；
  \item 其他 Stateful Platform Services。
\end{itemize}
“Platform Service”同样是 DAG Role，不是 Domain Directory。

\vspace{0.5em}\hrule\vspace{0.5em}
\subsection{10. Domain Ontology}
目录只表达 Domain Ownership。

\begin{lstlisting}[language=text]
gitops/platform/
├── applications/
├── foundation/
├── management/
├── operators/
├── networking/
├── observability/
├── messaging/
└── storage/
\end{lstlisting}
\vspace{0.5em}\hrule\vspace{0.5em}
\subsubsection{10.1 \_\_INLINE\_CODE\_0\_\_}
不是业务领域。

仅保存：

\begin{itemize}
  \item Application；
  \item Kustomization；
  \item Application DAG Overlay；
  \item Wave Metadata；
  \item Parent-level protection metadata。
\end{itemize}
\vspace{0.5em}\hrule\vspace{0.5em}
\subsubsection{10.2 \_\_INLINE\_CODE\_0\_\_}
Kubernetes 原生逻辑边界。

\vspace{0.5em}\hrule\vspace{0.5em}
\subsubsection{10.3 \_\_INLINE\_CODE\_0\_\_}
平台自身治理能力：

\begin{itemize}
  \item Argo CD Self Management；
  \item AppProjects；
  \item RBAC；
  \item OIDC；
  \item Sealed Secrets；
  \item ESO；
  \item cert-manager；
  \item Policy Engine。
\end{itemize}
AppProject Desired State 必须归属：

\begin{lstlisting}[language=text]
platform/management/projects/
\end{lstlisting}
Project Bootstrap Application 只是使用 \_\_INLINE\_CODE\_0\_\_ 去同步这些对象。

因此：

\begin{gfmBlockquote}
Bootstrap lifecycle ≠ AppProject domain ownership.
\end{gfmBlockquote}
\vspace{0.5em}\hrule\vspace{0.5em}
\subsubsection{10.4 \_\_INLINE\_CODE\_0\_\_}
领域扩展 API Controller：

\begin{itemize}
  \item Flink Operator；
  \item Spark Operator；
  \item 其他平台领域 Operator。
\end{itemize}
\vspace{0.5em}\hrule\vspace{0.5em}
\subsubsection{10.5 \_\_INLINE\_CODE\_0\_\_}
网络与流量能力：

\begin{itemize}
  \item Envoy Gateway；
  \item Gateway API；
  \item Gateway；
  \item HTTPRoute；
  \item NetworkPolicy-related resources。
\end{itemize}
\vspace{0.5em}\hrule\vspace{0.5em}
\subsubsection{10.6 \_\_INLINE\_CODE\_0\_\_}
统一遥测：

\begin{itemize}
  \item Kube-Prometheus-Stack；
  \item Alloy；
  \item OpenTelemetry Collector；
  \item Loki；
  \item Tempo；
  \item OTLP Transport。
\end{itemize}
\vspace{0.5em}\hrule\vspace{0.5em}
\subsubsection{10.7 \_\_INLINE\_CODE\_0\_\_}
异步事件总线：

\begin{itemize}
  \item Redpanda；
  \item Kafka。
\end{itemize}
\vspace{0.5em}\hrule\vspace{0.5em}
\subsubsection{10.8 \_\_INLINE\_CODE\_0\_\_}
数据持久化平台：

\begin{itemize}
  \item MinIO；
  \item 未来其他存储服务。
\end{itemize}
\vspace{0.5em}\hrule\vspace{0.5em}
\subsection{11. Leaf Application Contract}
每个可独立：

\begin{itemize}
  \item 失败；
  \item 暂停；
  \item 升级；
  \item 回滚；
  \item 观察；
  \item 恢复；
\end{itemize}
的平台能力 SHOULD 成为独立 Leaf Application。

Leaf Application 是：

\begin{gfmBlockquote}
Minimum Reconciliation and Failure Isolation Unit.
\end{gfmBlockquote}
禁止仅因为：

\begin{lstlisting}[language=text]
platform/observability/
\end{lstlisting}
存在子目录，就形成：

\begin{lstlisting}[language=text]
platform-app
 ↓
observability-app
 ↓
prometheus-app
\end{lstlisting}
这种目录镜像式 Application 链。

默认最大结构为：

\begin{lstlisting}[language=text]
External Root
      ↓
Platform Control
      ↓
Leaf Application
      ↓
Resources
\end{lstlisting}
任何额外 Application 层 MUST 由 ADR 证明其具有真实控制边界价值。

\vspace{0.5em}\hrule\vspace{0.5em}
\subsection{12. Application Health Gate Contract}
Sync Wave 只有结合 Health Gate 才构成真正的依赖图。

Atlas 的 DAG Contract 是：

\begin{lstlisting}[language=text]
Sync Wave
+
Child Automated Sync
+
Application Health
=
Deterministic Application DAG
\end{lstlisting}
三者缺一不可。

\vspace{0.5em}\hrule\vspace{0.5em}
\subsubsection{12.1 Application Health Capability}
Bootstrap Seed 与 \_\_INLINE\_CODE\_0\_\_ 必须注入同一份：

\begin{lstlisting}[language=text]
argoproj.io/Application
\end{lstlisting}
健康评估能力。

该 Capability MUST 来自同一份受版本控制资产。

禁止：

\begin{lstlisting}[language=text]
Bootstrap health semantics
≠
Steady-state health semantics
\end{lstlisting}
\vspace{0.5em}\hrule\vspace{0.5em}
\subsubsection{12.2 Child Automated Sync}
所有作为 DAG 前置条件的 Child Application：

\begin{gfmBlockquote}
MUST enable Automated Sync
\end{gfmBlockquote}
或者使用一个外部确定性触发机制。

禁止出现：

\begin{lstlisting}[language=text]
Wave -50 App
Healthy only after manual sync
\end{lstlisting}
却被用作：

\begin{lstlisting}[language=text]
Wave -30 prerequisite
\end{lstlisting}
否则 DAG 会永久阻塞。

\vspace{0.5em}\hrule\vspace{0.5em}
\subsubsection{12.3 Macro Health Gate}
Root 不需要理解 Platform 内部组件。

Root 只判断：

\begin{lstlisting}[language=text]
platform-control = Healthy
\end{lstlisting}
因此：

\begin{lstlisting}[language=text]
Foundation
+
Management
+
Operators
+
Controllers
+
Platform Services
        ↓
Converged
        ↓
Platform Control Healthy
        ↓
Workload Control may proceed
\end{lstlisting}
这将：

\begin{gfmBlockquote}
Platform Ready
\end{gfmBlockquote}
从人类经验判断转化为机器可判定状态。

\vspace{0.5em}\hrule\vspace{0.5em}
\subsection{13. Workload Control Model}
Workload Control 必须区分：

\begin{gfmBlockquote}
Control-plane Orchestrator
\end{gfmBlockquote}
与：

\begin{gfmBlockquote}
Tenant Payload Application.
\end{gfmBlockquote}
这是 Tier-1 / Tier-2 权限隔离的核心。

\vspace{0.5em}\hrule\vspace{0.5em}
\subsubsection{13.1 Workload Control Application}
Root 中：

\begin{lstlisting}[language=text]
workload-control-app.yaml
\end{lstlisting}
如果其职责是创建：

\begin{itemize}
  \item Application；
  \item ApplicationSet；
  \item 其他 Argo CD Control Objects；
\end{itemize}
则其本质属于：

\begin{gfmBlockquote}
Tier-1 Control Plane.
\end{gfmBlockquote}
因此：

\begin{lstlisting}[language=text]
project: platform-project
\end{lstlisting}
而不是：

\begin{lstlisting}[language=text]
project: workload-project
\end{lstlisting}
\vspace{0.5em}\hrule\vspace{0.5em}
\subsubsection{13.2 Tenant Leaf Applications}
由 Workload Control 生成或管理的最终业务 Application：

\begin{lstlisting}[language=text]
project: workload-project
\end{lstlisting}
其资源只能进入被允许的 Workload Namespace。

控制关系：

\begin{lstlisting}[language=text]
Workload Control
project: platform-project
          │
          ▼
Application / ApplicationSet
          │
          ▼
Tenant Leaf Application
project: workload-project
          │
          ▼
Developer Namespace
\end{lstlisting}
\vspace{0.5em}\hrule\vspace{0.5em}
\subsubsection{13.3 Why}
禁止让 \_\_INLINE\_CODE\_0\_\_ 为了创建 Argo CD Application/ApplicationSet 而获得：

\begin{itemize}
  \item Argo CD Namespace 写权限；
  \item Platform CR 写权限；
  \item 高危 Cluster Capability。
\end{itemize}
原则是：

\begin{gfmBlockquote}
WHO creates tenant Applications → Tier-1
WHAT tenant Applications may deploy → Tier-2
\end{gfmBlockquote}
\vspace{0.5em}\hrule\vspace{0.5em}
\subsection{14. Workload Directory}
建议标准：

\begin{lstlisting}[language=text]
gitops/workloads/
├── applications/
│   ├── base/
│   └── overlays/
│
└── apps/
    ├── clickstream-processor/
    ├── example-streaming-job/
    └── ...
\end{lstlisting}
其中：

\begin{lstlisting}[language=text]
applications/
\end{lstlisting}
保存 Workload Orchestration Metadata。

\begin{lstlisting}[language=text]
apps/
\end{lstlisting}
保存真正的业务 Desired State。

Phase 2B 可以将：

\begin{lstlisting}[language=text]
applications/
\end{lstlisting}
演进为 ApplicationSet Tenant Model。

\vspace{0.5em}\hrule\vspace{0.5em}
\subsection{15. Data Bootstrap Jobs}
数据初始化逻辑不得因为某个 Wave 到达就默认具有破坏性执行权。

安全初始化必须：

\begin{gfmBlockquote}
100\% Idempotent.
\end{gfmBlockquote}
例如：

\begin{itemize}
  \item create-if-not-exists；
  \item schema initialization；
  \item non-destructive seed。
\end{itemize}
\vspace{0.5em}\hrule\vspace{0.5em}
\subsubsection{15.1 Safe Initialization}
允许进入自动 DAG。

推荐作为 Workload Control Scope 中：

\begin{lstlisting}[language=text]
+10
\end{lstlisting}
级别的初始化单元。

\vspace{0.5em}\hrule\vspace{0.5em}
\subsubsection{15.2 Destructive Migration}
例如：

\begin{itemize}
  \item DROP；
  \item destructive schema rewrite；
  \item irreversible state conversion；
  \item mass data deletion；
\end{itemize}
不得依赖普通 AutoSync 自动执行。

必须转化为：

\begin{gfmBlockquote}
Human-gated Workflow.
\end{gfmBlockquote}
\vspace{0.5em}\hrule\vspace{0.5em}
\subsection{16. Sync Wave Scope}
Atlas 正式规定：

\begin{gfmBlockquote}
Sync Wave is parent synchronization-scope local.
\end{gfmBlockquote}
因此：

Root：

\begin{lstlisting}[language=text]
-110 Project Bootstrap
-100 Platform Control
   0 Workload Control
\end{lstlisting}
与 Platform Control：

\begin{lstlisting}[language=text]
-100 Foundation
 -90 Management
 -50 Operators
 -30 Controllers
 -10 Services
\end{lstlisting}
虽然存在相同数字：

\begin{lstlisting}[language=text]
-100
\end{lstlisting}
但它们不属于同一个全局调度空间。

禁止将所有 Wave 理解成：

\begin{gfmBlockquote}
Cluster-wide absolute sequence number.
\end{gfmBlockquote}
\vspace{0.5em}\hrule\vspace{0.5em}
\subsection{17. Environment Hydration}
Environment-specific state 与通用 Platform Domain 必须分离。

标准结构：

\begin{lstlisting}[language=text]
gitops/
├── root/
│   └── overlays/
│       ├── local-orbstack/
│       └── prod/
│
├── environments/
│   ├── local-orbstack/
│   └── prod/
│
└── platform/
    └── applications/
        └── overlays/
            ├── local-orbstack/
            └── prod/
\end{lstlisting}
\vspace{0.5em}\hrule\vspace{0.5em}
\subsubsection{17.1 Root Hydration}
Bootstrap 选择：

\begin{lstlisting}[language=text]
root/overlays/<environment>
\end{lstlisting}
作为 External Root Source。

\vspace{0.5em}\hrule\vspace{0.5em}
\subsubsection{17.2 Platform Hydration}
\_\_INLINE\_CODE\_0\_\_ 选择：

\begin{lstlisting}[language=text]
platform/applications/overlays/<environment>
\end{lstlisting}
作为 Platform DAG。

\vspace{0.5em}\hrule\vspace{0.5em}
\subsubsection{17.3 Environment-specific K8s State}
例如：

\begin{itemize}
  \item local StorageClass；
  \item Cloud-specific CSI；
  \item environment-specific Gateway address；
  \item substrate-specific configuration；
\end{itemize}
必须位于：

\begin{lstlisting}[language=text]
gitops/environments/<environment>/
\end{lstlisting}
或由对应 Overlay 显式引用。

不得污染通用 Platform Domain。

\vspace{0.5em}\hrule\vspace{0.5em}
\subsection{18. Bootstrap Adoption Contract}
Bootstrap 恢复必须满足：

\begin{lstlisting}[language=text]
Object Identity Parity
+
Version Parity
+
Health Capability Parity
+
CRD Compatibility
=
Bootstrap Adoption Contract
\end{lstlisting}
\vspace{0.5em}\hrule\vspace{0.5em}
\subsubsection{18.1 \_\_INLINE\_CODE\_0\_\_}
Bootstrap 与 GitOps Self-management 必须共同消费：

\begin{lstlisting}[language=text]
versions.lock
\end{lstlisting}
禁止：

\begin{lstlisting}[language=text]
bootstrap Argo version
≠
argocd-self desired version
\end{lstlisting}
所导致的隐式：

\begin{itemize}
  \item downgrade；
  \item upgrade；
  \item CRD incompatibility。
\end{itemize}
\vspace{0.5em}\hrule\vspace{0.5em}
\subsubsection{18.2 Recovery Seed}
Break-Glass Seed 的目标是：

\begin{gfmBlockquote}
Minimal but adoption-compatible.
\end{gfmBlockquote}
不是重新建立另一套独立控制平面。

\vspace{0.5em}\hrule\vspace{0.5em}
\subsection{19. Dual Deletion Protection}
Atlas 实施两层独立保护。

\vspace{0.5em}\hrule\vspace{0.5em}
\subsubsection{19.1 Parent-level Prune Protection}
所有由 Parent Application 管理的 Tier-0 / Tier-1 Child Application CR 应具有：

\begin{lstlisting}[language=text]
argocd.argoproj.io/sync-options: Prune=confirm
\end{lstlisting}
保护链：

\begin{lstlisting}[language=text]
Git accidentally removes child Application
                 ↓
Parent attempts prune
                 ↓
Prune suspended
                 ↓
Human confirmation
\end{lstlisting}
\vspace{0.5em}\hrule\vspace{0.5em}
\subsubsection{19.2 No Cascading Finalizer}
核心 Tier-0 / Tier-1 Application 默认不得携带级联资源删除 finalizer。

因此：

\begin{lstlisting}[language=text]
delete Application
        ↓
Application CR deleted
        ↓
Managed resources survive
\end{lstlisting}
\vspace{0.5em}\hrule\vspace{0.5em}
\subsubsection{19.3 Protection Boundary}
上述保护主要针对：

\begin{gfmBlockquote}
Application Control Graph.
\end{gfmBlockquote}
并不意味着：

\begin{lstlisting}[language=text]
StatefulSet
PVC
Kafka Cluster CR
Database
Namespace
\end{lstlisting}
自动获得相同保护。

对于高价值 Leaf Resources，必须根据风险显式配置：

\begin{itemize}
  \item resource-level \_\_INLINE\_CODE\_0\_\_；
  \item \_\_INLINE\_CODE\_0\_\_；
  \item Admission Policy；
  \item Domain-specific retention mechanism。
\end{itemize}
\vspace{0.5em}\hrule\vspace{0.5em}
\subsection{20. Secret Dependency Contract}
Phase 1 使用 Sealed Secrets 时：

\begin{lstlisting}[language=text]
Sealed Secrets Controller
          ↓
Controller Healthy
          ↓
SealedSecret Reconciliation
          ↓
Secret Exists
          ↓
Consumer Application
\end{lstlisting}
Consumer 不得在 Secret Capability Ready 前进入强依赖路径。

Phase 2A 切换 ESO 时同样遵守：

\begin{lstlisting}[language=text]
ESO
 ↓
External Secret Materialization
 ↓
Consumer
\end{lstlisting}
\vspace{0.5em}\hrule\vspace{0.5em}
\subsection{21. Emergency Recovery \& Blast Radius Freeze}
Atlas 使用：

\begin{gfmBlockquote}
Freeze Outside-In, Resume Inside-Out.
\end{gfmBlockquote}
\vspace{0.5em}\hrule\vspace{0.5em}
\subsubsection{21.1 Detect}
确认：

\begin{itemize}
  \item 错误 Revision；
  \item 受影响 Application；
  \item 数据面状态；
  \item 爆炸半径；
  \item 是否仍有危险 AutoSync。
\end{itemize}
\vspace{0.5em}\hrule\vspace{0.5em}
\subsubsection{21.2 Freeze Outside-In}
首先冻结：

\begin{lstlisting}[language=text]
External Root Anchor
\end{lstlisting}
然后依次关闭受影响 Control Application 的 AutoSync。

重点包括：

\begin{itemize}
  \item Platform Control；
  \item argocd-self；
  \item 受影响 Leaf Applications；
  \item Workload Control。
\end{itemize}
目标不是删除资源，而是：

\begin{gfmBlockquote}
Stop propagation.
\end{gfmBlockquote}
\vspace{0.5em}\hrule\vspace{0.5em}
\subsubsection{21.3 Snapshot}
备份：

\begin{itemize}
  \item Application 状态；
  \item AppProject 状态；
  \item 当前 Revision；
  \item Controller Configuration；
  \item 关键 Event / Audit Evidence。
\end{itemize}
\vspace{0.5em}\hrule\vspace{0.5em}
\subsubsection{21.4 Restore Seed}
使用 Bootstrap 恢复满足 Adoption Contract 的 Argo CD Seed。

\vspace{0.5em}\hrule\vspace{0.5em}
\subsubsection{21.5 Repair Git}
Git 必须回到：

\begin{gfmBlockquote}
Known-Good Revision.
\end{gfmBlockquote}
不得依赖长期 Live Edit 作为恢复终态。

\vspace{0.5em}\hrule\vspace{0.5em}
\subsubsection{21.6 Resume Inside-Out}
推荐顺序：

\begin{lstlisting}[language=text]
Argo Seed
   ↓
argocd-self
   ↓
Management Capability
   ↓
Platform Capability DAG
   ↓
Platform Control Healthy
   ↓
Workload Control
   ↓
Tenant Workloads
   ↓
External Root normal reconciliation
\end{lstlisting}
必须逐级验证 Health。

\vspace{0.5em}\hrule\vspace{0.5em}
\subsection{22. ApplicationSet Boundary}
ApplicationSet 是 Atlas Phase 2B Tenant Platform 的核心候选机制，但不作为 v1.0.3 平台 Bootstrap Correctness 的必要基础。

当前 Frozen Baseline 的平台启动正确性依赖：

\begin{lstlisting}[language=text]
Application
+
App-of-Apps
+
Sync Wave
+
Application Health
\end{lstlisting}
ApplicationSet 的主要目标为：

\begin{itemize}
  \item Tenant Expansion；
  \item Multi-cluster Workload Generation；
  \item Developer Self-Service；
  \item Fleet-style Application Generation。
\end{itemize}
ApplicationSet 不应取代 Tier-0 Root Trust Model。

\vspace{0.5em}\hrule\vspace{0.5em}
\subsection{23. AI \& Human Control Boundary}
\_\_INLINE\_CODE\_0\_\_：

\begin{lstlisting}[language=text]
READ      allowed
ANALYZE   allowed
PROPOSE   allowed
GENERATE  allowed
MERGE     human-gated
APPLY     human-gated
\end{lstlisting}
Tier-1：

允许未来在 Policy Guardrail 成熟后开放更高程度 Agent Autonomy。

Tier-2：

是最适合 AI Agent 自治执行的主要区域。

治理演进顺序必须是：

\begin{lstlisting}[language=text]
Guardrail
   ↓
Auditability
   ↓
Policy Enforcement
   ↓
Agent Autonomy
\end{lstlisting}
而不是先开放自治，再补安全机制。

\vspace{0.5em}\hrule\vspace{0.5em}
\subsection{24. Phase 2 Governance Roadmap}
\vspace{0.5em}\hrule\vspace{0.5em}
\subsubsection{Phase 2A — Safety Foundation}
引入：

\begin{itemize}
  \item Kyverno / equivalent policy engine；
  \item NetworkPolicy baseline；
  \item ESO；
  \item Root Anchor admission protection；
  \item stronger audit controls。
\end{itemize}
\vspace{0.5em}\hrule\vspace{0.5em}
\subsubsection{Phase 2B — Tenant Platform}
建设：

\begin{itemize}
  \item ApplicationSet；
  \item Tenant Model；
  \item Namespace lifecycle；
  \item workload-project policy；
  \item multi-tenant boundaries。
\end{itemize}
\vspace{0.5em}\hrule\vspace{0.5em}
\subsubsection{Phase 2C — Developer Experience}
建设：

\begin{itemize}
  \item Backstage；
  \item Golden Path；
  \item Templates；
  \item Platform API；
  \item developer self-service。
\end{itemize}
\vspace{0.5em}\hrule\vspace{0.5em}
\subsubsection{Phase 2D — Agent Automation}
在 Guardrail 完成后，开放：

\begin{itemize}
  \item Agent-generated PR；
  \item Policy-verified changes；
  \item limited auto-merge；
  \item bounded autonomous remediation；
  \item declarative platform operation。
\end{itemize}
Agent 自治权不得跨越 Tier-0 Human Judgment Boundary。

\vspace{0.5em}\hrule\vspace{0.5em}
\subsection{25. GitOps Control Plane Invariants}
Atlas GitOps Control Plane v1.0.3 冻结以下领域不变量。

\subsubsection{GITOPS-01}
External Root Anchor MUST remain outside parent Application reconciliation.

\subsubsection{GITOPS-02}
Canonical AppProjects MUST be:

\begin{lstlisting}[language=text]
atlas-bootstrap
platform-project
workload-project
\end{lstlisting}
\subsubsection{GITOPS-03}
Root MUST contain only macro trust/readiness control Applications.

\subsubsection{GITOPS-04}
Platform components MUST be orchestrated by Platform Control rather than enumerated directly under Root.

\subsubsection{GITOPS-05}
Directory hierarchy MUST NOT be mechanically mirrored as Application hierarchy.

\subsubsection{GITOPS-06}
Leaf Application SHOULD be the minimum failure-isolation unit.

\subsubsection{GITOPS-07}
Sync Wave MUST be interpreted within its parent synchronization scope.

\subsubsection{GITOPS-08}
Any Child Application acting as a DAG prerequisite MUST have deterministic synchronization.

\subsubsection{GITOPS-09}
Bootstrap and argocd-self MUST share Application Health semantics.

\subsubsection{GITOPS-10}
Bootstrap and argocd-self MUST share \_\_INLINE\_CODE\_0\_\_.

\subsubsection{GITOPS-11}
Workload orchestrators creating Argo CD control objects MUST operate in Tier-1.

\subsubsection{GITOPS-12}
Tenant payload Applications MUST operate in \_\_INLINE\_CODE\_0\_\_.

\subsubsection{GITOPS-13}
Tier-0 / Tier-1 control graph deletion MUST NOT implicitly destroy descendant runtime resources.

\subsubsection{GITOPS-14}
Destructive data mutation MUST NOT be triggered merely by automatic Wave progression.

\vspace{0.5em}\hrule\vspace{0.5em}
\subsection{26. Final Frozen Control Graph}
Atlas v1.0.3 的最终规范拓扑为：

\begin{lstlisting}[language=text]
                         Bootstrap
                             │
                 ┌───────────┴───────────┐
                 ▼                       ▼
             Argo Seed             atlas-bootstrap
                                        │
                                        ▼
                             External Root Anchor
                                        │
                         Root Macro Sync Scope
                                        │
                 ┌──────────────────────┼─────────────────────┐
                 │                      │                     │
                 ▼                      ▼                     ▼
       -110 Project Bootstrap   -100 Platform Control    0 Workload Control
          atlas-bootstrap          platform-project        platform-project
                 │                      │                     │
        ┌────────┴────────┐             │                     │
        ▼                 ▼             ▼                     ▼
platform-project   workload-project   Platform DAG       Workload Orchestration
                                        │                     │
                           Platform Sync Scope                  │
                                        │                     ▼
                    ┌───────────────────┼──────────────┐  Tenant Applications
                    │                   │              │   workload-project
                    ▼                   ▼              ▼
               Foundation          Management       Operators
                  -100                -90             -50
                                                           │
                                                           ▼
                                                      Controllers
                                                         -30
                                                           │
                                                           ▼
                                                   Platform Services
                                                         -10
                                                           │
                                                           ▼
                                                     Platform Healthy
                                                           │
                                                           ▼
                                                     Workloads Enabled
\end{lstlisting}
这张图是 Atlas GitOps Control Plane v1.0.3 的架构冻结点。

未来：

\begin{itemize}
  \item 增加 Redpanda；
  \item 增加 ClickHouse；
  \item 增加 Spark；
  \item 增加 Loki；
  \item 增加 Tempo；
  \item 增加新的 Operator；
\end{itemize}
均应优先修改：

\begin{lstlisting}[language=text]
gitops/platform/
\end{lstlisting}
而不是：

\begin{lstlisting}[language=text]
gitops/root/
\end{lstlisting}
只有当：

\begin{itemize}
  \item Trust Domain 改变；
  \item Control Ownership 改变；
  \item Macro Readiness Boundary 改变；
  \item Bootstrap Chain 改变；
\end{itemize}
时，才允许修改 Tier-0 Root Architecture。

\vspace{0.5em}\hrule\vspace{0.5em}
\subsection{27. Architecture Freeze Statement}
Atlas GitOps v1.0.3 正式冻结以下模式：

\begin{gfmBlockquote}
\textbf{External Root Anchor

- Minimal Root Macro DAG
- Flat Platform Capability DAG
- Tier-1 Workload Orchestration
- Tier-2 Workload Execution}
\end{gfmBlockquote}
其核心目标不是创建最复杂的 GitOps 拓扑，而是创建：

\begin{gfmBlockquote}
\textbf{最容易证明、最容易恢复、最容易审计、最不容易被人类或 AI Agent 误操作的控制图。}
\end{gfmBlockquote}
从本版本开始，任何将 Root 重新扩张为 Platform Component Catalog，或仅为了匹配目录层次而引入深层 App-of-Apps 的实现，应默认视为对 Frozen Baseline 的偏离，并进入正式 Architecture Review。

\end{document}